\scanpage{114}
\mspage{59}
% archive sheet number: 113

dans $\underline{\mathrm{Fibsc}}\,S$ (comme aussi dans $\underline{\mathrm{Fib}}\,S$ et $\underline{\mathrm{Liss}}\,S$) avec la $W$-équivalence par fibres, et d'autre part $\Phi(f)_s$ n'est autre que la flèche \unsure{précédente|précédent} $ {}_{s\backslash}X\mathbin{\unsure{-|\to}}{}_{s\backslash}Y$. Donc il est évident que \unsure{if|$f$|$\updownarrow$}

\[
 f\in\underset{\substack{\begin{tikzpicture}[baseline=0pt]\draw (0,.065)--(0,0)--(.65,0)--(.65,.065);\end{tikzpicture}\\W\text{-équiv.}\\ \text{droite}}}{W_S^d}\iff
 \Phi(f)\in\underset{\substack{\begin{tikzpicture}[baseline=0pt]\draw (0,.065)--(0,0)--(.65,0)--(.65,.065);\end{tikzpicture}\\W\text{-équiv.}\\ \text{par fibres}}}{W_S^f}
\]

b) Que $X\xrightarrow{i_X}\underset{\substack{\begin{tikzpicture}[baseline=0pt]\draw (0,.065)--(0,0)--(1.05,0)--(1.05,.065);\end{tikzpicture}\\=\Phi(X)}}{\lambda\Phi(X)}$ est dans $W_S^d$, pour tout $X$ dans $\mathrm{Cat}_{/S}$. \unsure{Cela|Celà} signifie donc que $i_X$ induit

\[
 {}_{s\backslash}X\longrightarrow{}_{s\backslash}\Phi(X)
\]

qui est une $W$-équivalence. Ce serait o.k. s'il était clair que $X\longmapsto\Phi(X)$ commute \deleted{à} tout chg de base \unsure{proposé|propre} sur $S$ (donc en particulier à $ {}_{s\backslash}S\longrightarrow S$), puisqu'on sait que $X\mathbin{\unsure{-|\to}}\Phi(X)$ est une \inserted{$W$-}équivalence \deleted{\unsure{$W$}}. C'est sans doute trop demander, mais du moins que ça commute au chgt de \deleted{base} qui est un co-isom. local (telles les $ {}_{s\backslash}S\longrightarrow S$).

\deleted{Or} Pour un tel chgt. de base $S'\xrightarrow{p}S$, vérifier que \deleted{\unsure{les}} \unsure{la|les} \unsure{fibre}\deleted{s} de $\Phi(X')$ \inserted{en $s'$} \deleted{\unsure{sont celles}} \inserted{est} \deleted{de} isom. à celle de $\Phi(X)$ en $s\inserted{=p(s')}$ \unsure{n'en} coûte rien. Or cette fibre est $ {}_{s'\backslash}X'$, déduite de $X'$ sur $S'$ \deleted{$\unsure{{}_{s'\backslash}\Phi(X')=\Phi}$} par le chgt. de base $ {}_{s'\backslash}S'\mathbin{\unsure{-|\to}} S'$, \inserted{donc déduit de $X$ par $ {}_{s'\backslash}S'\longrightarrow\unsure{S}$,} Or $ {}_{s'\backslash}S'\longrightarrow{}_{s\backslash}S$ est un isom par hyp.

